\documentclass{article}
\usepackage{amsmath, amssymb, amsthm}
\usepackage{hyperref}
\usepackage{geometry}
\geometry{margin=1in}

\title{Erd\H{o}s Problem \#342}
\date{Last updated: 2025-08-31}
\author{Source: erdosproblems.com}

\begin{document}
\maketitle

\noindent\textbf{Status:} OPEN \quad \textbf{No prize}

\noindent\textbf{Tags:} number theory

\medskip
\noindent\textbf{Problem Statement:}

With $a_1=1$ and $a_2=2$ let $a_{n+1}$ for $n\geq 2$ be the least integer $>a_n$ which can be expressed uniquely as $a_i+a_j$ for $i<j\leq n$.

What can be said about this sequence? Do infinitely many pairs $a,a+2$ occur? Does this sequence eventually have periodic differences? Is the density $0$?

\bigskip
\noindent\textbf{Remarks:}

A problem of Ulam. The sequence is\[1,2,3,4,6,8,11,13,16,18,26,28,\ldots\]at OEIS A002858.

See also Problem 7 of Green's open problems list.

This is problem C4 in Guy's collection \cite{Gu04}.

\bigskip
\noindent\textbf{References:}

[Gu04] Guy, Richard K., Unsolved problems in number theory. (2004), xviii+437.

\bigskip
\noindent\small{Source: \url{https://www.erdosproblems.com/342}}
\end{document}
